% Section 1 -- Introduction.
% Ported from v2 (legacy/v2_from_rp_kisan/experiments/paper_build/content.js) and rewritten
% per paper/WRITING_GUIDE.md. Citations still use the v2 reference list; the literature
% session (R1) verifies them in literature/ before submission.
\section{Introduction}
\label{sec:intro}

\draftfor{Drafted by Claude from the v2 text. One author rewrites this section in their own
words and adds what only the team knows: why we started with Patiala, what FPO staff told us,
and what went wrong in v2. Delete this box when done.}

A Farmer Producer Organization (FPO) is a registered body of farmers that can pool, market
and sell its members' produce~\cite{sfac}. Selling the produce after harvest is where many
FPOs struggle. For each member lot, the FPO has to choose a mandi, a day and a quantity, and
these choices depend on each other. The mandi with the best price may be the farthest from the
village. A later selling date may pay more, but grain kept in storage can heat up or gain
moisture. When many members deliver in the same week, a single mandi can receive more than it
can weigh and unload in a day.

Losses after harvest are large in India. A 2024 ICRIER policy brief puts them at about
Rs~1.53 trillion of food a year, using a national NABCONS survey from 2020--2022, and it names
storage as one of the causes~\cite{icrier}.
\authornote{Literature session: re-check this figure against the brief before submission.}

Two public sources give an FPO most of the price information it needs. AGMARKNET publishes the
daily arrivals and the minimum, maximum and modal prices reported by each mandi~\cite{agmarknet}.
For paddy and wheat, the government also announces a minimum support price (MSP) before each
marketing season~\cite{msp1,msp2}. For Paddy (Common) in Punjab, the mandi prices hardly
differ from it. Across the \kn{mandis} mandis we studied, at least \kn{msp_share_min}\% of the reported prices
in every season from \kn{season_first} to \kn{season_last} were exactly the Grade-A MSP. Basmati
is not covered by the MSP, and wheat that farmers hold after the procurement window is sold at
market prices, so for these crops the timing of a sale can change what a farmer earns.
\authornote{Cite a source for the wheat procurement window and basmati pricing.}

KisanNiti is the decision-support system we are building for FPO staff in the Patiala area.
It reads AGMARKNET prices, the FPO's records of member lots and readings from a grain-storage
sensor, and proposes a selling plan that a manager reviews before farmers are told on WhatsApp
(Fig.~\ref{fig:architecture}).
\authornote{Contributions paragraph: write it when the v3 results exist. Name each result
with its section, e.g.\ the paddy forecasting comparison (Sect.~\ref{sec:results-msp}), the
wheat and basmati comparison, the storage-risk rule, and the fair allocation model. Avoid a
``First \ldots\ Second \ldots\ Third'' list.}

Section~\ref{sec:related} reviews related work. Section~\ref{sec:system} describes the system
and our methods, Sect.~\ref{sec:results} reports the experiments, and Sect.~\ref{sec:conclusion}
states the limits of this study and the next steps.
